% SOURCE_AUTHORITY: sga5_fr_workpass.tex, lines 15246--15246 (LF-normalized unit).
% SOURCE_SHA256: 791F4EFFC5E02832D5D77ED03518C8156D6F07E4C8238B03545DB93D883FBB28
Para demostrar (ii), se comienza observando que el caso en que $P_0$ es libre es evidente (sin hipótesis alguna sobre el ideal $I$). Se pasa al caso general en que $P_0$ es proyectivo introduciendo un $A_0$-módulo libre $L_0$ del que $P_0$ es sumando directo y un $A$-módulo libre $L$ tal que $T(L)\simeq L_0$. El módulo $P_0$ se identifica con la imagen de un proyector $e_0:L_0\to L_0$, y hay que ver que existe un proyector $e:L\to L$ que levanta $e_0$, lo cual resulta de Bourbaki, Alg. comm., cap.~III, §4, n\textsuperscript{o}~6, lema 2, aplicado a la sub-$A$-álgebra $B$ de $\End(L)$ engendrada por un endomorfismo de $L$ que levanta $e_0$ (cuya existencia está garantizada por (i)) y al ideal (nilpotente) dado por la intersección de $B$ con el núcleo del homomorfismo canónico $\End(L)\to\End(L_0)$. Esto demuestra la existencia de $P$ en (ii). Sea además $P'$ como enunciado en (ii); entonces el isomorfismo compuesto $T(P)\simeq P_0\simeq T(P')$ procede, en virtud de (i), de un homomorfismo $f:P\to P'$. Como $P$ y $P'$ son planos sobre $A$ y $T(f)$ es un isomorfismo, se concluye que $\operatorname{Gr}_I(f)$ es un isomorfismo, y por tanto $f$ es un isomorfismo (pues $I$ es nilpotente), de donde se sigue (ii).
